% id: 135-02
% book: 베이직쎈 확률과 통계 · 07 통계적 추정
% section: 유형 083 모평균의 추정; 모표준편차가 주어진 경우
% figure: tikz:fig-135-02
% answer: 
% answer_source: 
% variant_level: 0 (원본 전사)
% 이 파일은 build.mjs 가 items.json 에서 만든다. 고칠 때는 items.json 을 고친다.
\smash{\raisebox{1.5mm}{\dmkichul{베이직쎈 확률과 통계 135쪽 02번}}}
\noindent\begin{minipage}[t]{0.58\linewidth}\vspace{0pt}\raggedright 정규분포 $\mathrm{N}(m{,}\mkern3mu\mathopen{}10^2)$을 따르는 모집단에서 크기가 225인 표본을 임의추출하여 구한 표본평균이 42일 때, 모평균 $m$에 대한 신뢰도 $99\mkern3mu \%$의 신뢰구간이 $42-k\le m\le 42+k$이다. 이때 위의 표준정규분포표를 이용하여 상수 $k$의 값을 구하면?\end{minipage}\hfill\begin{minipage}[t]{0.4\linewidth}\vspace{0pt}\centering \resizebox{\ifdim\width>\linewidth\linewidth\else\width\fi}{!}{% 135-02(인쇄 135쪽) — 135-02 오른쪽에 놓인 표준정규분포표($z$ / $\mathrm{P}(0\le Z\le z)$).
% 스캔 표라 크롭하지 않고 tabular 로 다시 조판했다(칸 값은 원문 그대로 · 원문에 없는 행은 넣지 않는다).
{\setlength{\tabcolsep}{5pt}\renewcommand{\arraystretch}{1.25}%
\begin{tabular}{|c|c|}
\hline
$z$ & $\mathrm{P}(0\le Z\le z)$ \\
\hline
$1.65$ & $0.450$ \\
\hline
$1.96$ & $0.475$ \\
\hline
$2.58$ & $0.495$ \\
\hline
\end{tabular}}
}\end{minipage}\par
\dmchoicesauto{}{$1.71$}{$1.72$}{$1.73$}{$1.74$}{$1.75$}
